\scanpage{5}
\mspage{- 110 -}
% archive sheet number: 4
% lines 1–2
\margin{\handwritten{abusif}}$\dim(X\handwritten{\times}_k Y)=\dim(X)+\dim(Y)$ ; \uline{en particulier, pour tout corps} $K$, \uline{extesion de} $k$, $\dim(X\handwritten{\otimes}_k K)=\dim(X)$.

% lines 3–5
Si $(U_i)$ et $(V_j)$ sont des recouvrements ouverts affines de $X$ et $Y$, $(U_i\handwritten{\times}_k V_j)$ est un recouvrement ouvert affine de $X{}_kY$, donc on peut se borner au cas où $X$ et $Y$ sont affines (n°1,cor.1 de la prop.2).

% lines 6–11
D'autre part, si \deleted{$\unsure{X}_i$} $X_i$ (resp. $Y_j$) sont les composantes irréductibles de $X$ (resp. $Y$), $X\handwritten{\times}_kY$ est réunion des $X_i\handwritten{\times}_kY_j$, qui sont ferm\handwritten{és} (\S{} 2,n°2,prop.5), donc (n°1,prop.1 (ii)) on peut supposer $X$ et $Y$ i\handwritten{r}réductibles ; \deleted{\illegible{1}} enfin, comme $(X\handwritten{\times}_kY)_{\mathrm{red}}$ et $(X_{\mathrm{red}}\handwritten{\times}_kY_{\mathrm{red}})_{\mathrm{red}}$ sont identiques (\S{} 3,n°4,cor. de la prop.11), on peut supposer $X$ e\handwritten{t} $Y$ réduits.

% lines 11–17
Soient donc $X=\mathrm{Spec}(A)$, $Y=\mathrm{Spec}(B)$, où $A$ et $B$ sont des algèbres intègres de type fini sur $k$. Si $m=\dim(X)$, $n=\dim(Y)$, il résulte du lemme de normalisation (Bourbaki,Alg.comm.\qquad) que $A$ (re\handwritten{sp.} $B$) est \deleted{\illegible{1}} un module de type fini sur un sous-anneau $A_1=\allowbreak =k\handwritten{[}a_1,\dots,a_m\handwritten{]}$ (resp. $B_1=k\handwritten{[}b_1,\dots,b_n\handwritten{]}$) où les $a_i$ (resp. $b_j$) sont algébriquement indépendants sur $k$ ; $C=A\handwritten{\otimes}_kB$ est donc un module de type fini sur $C_1=A_1\handwritten{\otimes}_kB_1$.
% slanted
\margin{\handwritten{Autre démonstration : les $\unsure{\kappa|K}(z)$ $(z\in X\times_k Y)$ sont toutes les ext. composées des ext. $\kappa(x)$ et $\kappa(y)$ $(x\in X,y\in Y)$, d'où aussitôt le résultat vu la théorie des corps}}

% lines 17–23
Pour tout idéal premier minimal $\handwritten{\underline}{p}$ \deleted{\unsure{dans}} dans $C$\handwritten{,} on voit par suite que le degré de transcendance \inserted{sur $k$} du corps des fracti\handwritten{ons} de $C/\handwritten{\underline}{p}$ est au plus $m+n$ ; \deleted{\illegible{1}} d'autre part, en vertu du th. de Cohen-Seidenberg,il existe un idéal premier $\handwritten{\underline}{q}$ de $C$ tel que $\handwritten{\underline}{q}\handwritten{\cap} C_1=\handwritten{\{}0\handwritten{\}}$, de il y a un des idéaux premiers minimaux \deleted{\unsure{dans}} dans \deleted{\illegible{1}} $C$ tel que le corps des fractions de $C/\handwritten{\underline}{p}$ soit de degré fini sur le corps des fractions de $C_1$, et par suite de degré de transcendance $m+n$ sur $k$.

% lines 24–28
\uline{Proposition 8 .- Soient} \inserted{\handwritten{$X$,}} $Y$\handwritten{,} \uline{\deleted{un} \inserted{\handwritten{deux}} préschéma\handwritten{s} localement noethérien\handwritten{s}}, $f:X\handwritten{\to}Y$ \uline{un morphisme \deleted{de type fini}. Soient} $y\handwritten{\in}Y$, \uline{et soit $x$ le point nérique d'une composante irréductible de} $f^{-1}(y)$, \uline{considéré comme préschéma sur} $\handwritten{\underline{\kappa}}(y)$ (\S{} 2,n°7,prop.17). \uline{Alors} $\operatorname{codim}_X\handwritten{\overline}{\handwritten{\{}x\handwritten{\}}}\handwritten{\leq}\operatorname{codim}_Y\handwritten{\overline}{\handwritten{\{}y\handwritten{\}}}$ (ou, ce qui revient au même, $\dim\handwritten{\underline}{O}_x\handwritten{\leq}\dim\handwritten{\underline}{O}_y$).
% slanted
\margin{\handwritten{\unsure{généraliser|Simplifier} en la formule $\dim\underline O_x\leq\dim\underline O_y+\dim\underline O_x$ (sur la fibre}}

% lines 29
La question est locale en $x$ et en $y$, puisque toute composante
% slanted
\margin{\handwritten{\deleted{\unsure{dire que la} \unsure{différence} \unsure{codim} dans $f^{-1}(y)$ \unsure{est majorée|est minorée} par la codim de \unsure{son adh.|plus un}, \unsure{ds|si} \unsure{son|non} $X$}}}
\margin{\handwritten{donner le cas d'égalité par \uline{platitude}}}
